This work investigated a technique designed to achieve two main objectives: re-synthesizing real musical instrument sounds and assessing the generalization ability of a model trained on randomized FM synthesizer parameters. The proposed Convolutional Neural Network (CNN) performed effectively in both tasks, reproducing instrumental timbres with high fidelity and generalizing well to unseen parameter configurations.

Comparisons with the NSynth dataset confirmed the model’s ability to approximate realistic spectral characteristics across different instrument families. The stability of the normalized Log-Mel metric further indicates that the model does not overfit specific timbral structures, supporting its robustness and flexibility. Overall, the results validate the initial hypothesis that CNNs can efficiently extract relevant spectral and temporal representations from audio signals and that randomized dataset generation can successfully guide the learning of FM synthesis parameters.

In addition to the experimental results, this work reinforces the idea that FM synthesis remains relevant for contemporary music technology. Although the technique is traditionally controlled through an empirical trial-and-error process, the use of machine learning can reduce the effort required to explore the parameter space and can make the synthesis workflow more accessible for musicians and composers.

Despite these positive outcomes, several limitations suggest directions for further research. The dataset was focused on musical instruments, which restricts broader generalization and motivates experiments with more diverse sound sources, such as environmental and synthetic noises. The exclusive use of a linear ADSR envelope also limited the diversity of dynamic profiles, indicating the need for richer envelope variations. Furthermore, only two FM synthesis algorithms were explored, and satisfactory performance was achieved with a simple configuration; investigating more complex operator structures may increase synthesis versatility.

Future work should explore alternative neural architectures, such as autoencoders, and comparative studies with RNNs and TCNs to improve feature representation and generalization. Integrating CNN-based extraction with other strategies, such as STFT-based regression inputs or perceptually informed loss functions, may further enhance performance. AutoML techniques could also optimize the regression stage. A more expressive FM synthesizer implementation, potentially a professional-grade one, is also desired to provide a more realistic evaluation framework. After such an improvement, a subjective analysis with human listeners could complement the objective metrics and validate perceptual quality. This direction would further advance neural sound synthesis for both academic and creative applications.
